% ===== latex-recomposed:document =====
% flm: style file flmrecomposed.6aba9645.sty
\flmHeading[toc=true]
\section{Tables of cells}{\flmLabels{\label{ref1}}}

A table between paragraphs, the first example of the Python product's documentation:

\flmCellsTable{cc}{X[-1]X[-1]}{hspan=minimal,
  cell{1}{1}={}{m, font={\flmCellsHeaderFont}},
  cell{1}{2}={}{m, font={\flmCellsHeaderFont}},
  hline{2}={1}{.4pt,solid},
  hline{2}={2}{.4pt,solid}}{\toprule
Header 1&Header 2\\
Row 1, Col 1&Row 1, Col 2\\
Row 2, Col 1&Row 2, Col 2\\
\bottomrule
}

A header spanning two columns, then data placed from the second column on:

\flmCellsTable{ccccc}{X[-1]X[-1]X[-1]X[-1]X[-1]}{hspan=minimal,
  cell{1}{3}={c=2}{l, font={\flmCellsHeaderFont}},
  hline{2}={3-4}{.4pt,solid}}{\toprule
&&Hi!&&\\
&A&B&C&D\\
&E&F&G&H\\
\bottomrule
}

Merged cells from a mapping:

\flmCellsTable{cccccc}{X[-1]X[-1]X[-1]X[-1]X[-1]X[-1]}{hspan=minimal,
  cell{1}{4}={r=2}{m},
  cell{1}{6}={r=2}{m},
  cell{3}{4}={r=2}{m},
  cell{3}{6}={r=2}{m}}{\toprule
&&&A&&AA\\
&&&&&\\
&&&B&&BB\\
&&&&&\\
\bottomrule
}

Styles per column, a mapping permuting the columns, and a cell overriding its entry:

\flmCellsTable{cccc}{X[-1]X[-1]X[-1]X[-1]}{hspan=minimal,
  cell{1}{1}={}{l, font={\flmCellsHeaderFont}},
  cell{1}{2}={}{c, font={\flmCellsHeaderFont}},
  cell{1}{3}={}{c, font={\flmCellsHeaderFont}},
  cell{3}{1}={}{m, font={\flmCellsHeaderFont}},
  cell{3}{2}={}{m, font={\flmCellsHeaderFont}},
  cell{3}{3}={c=2}{m, font={\flmCellsHeaderFont}},
  hline{2}={1}{.4pt,solid},
  hline{2}={2}{.4pt,solid},
  hline{2}={3}{.4pt,solid}}{\toprule
One&Two&Three&\\
C&B&A&\\
A&B&\flmCellsSet{green}C&\\
\bottomrule
}

Every token, with names for the rows and the columns:

\flmCellsTable{ccc}{X[-1]X[-1]X[-1]}{hspan=minimal,
  cell{1}{1}={}{m, font={\flmCellsHeaderFont}},
  cell{1}{2}={}{m, font={\flmCellsHeaderFont}},
  cell{1}{3}={}{m, font={\flmCellsHeaderFont}},
  cell{2}{1}={}{m, font={\flmCellsHeaderFont}},
  cell{2}{2}={}{l},
  cell{2}{3}={}{r},
  cell{3}{1}={}{m, font={\flmCellsHeaderFont}},
  cell{3}{2}={}{c},
  cell{3}{3}={}{r},
  cell{4}{1}={c=2}{m},
  cell{4}{3}={}{c, font={\flmCellsHeaderFont}},
  cell{5}{2}={c=2}{m},
  hline{2}={1}{.4pt,solid},
  hline{2}={2}{.4pt,solid},
  hline{2}={3}{.4pt,solid},
  vline{2}={3}{.4pt,solid},
  vline{3}={3}{.4pt,solid},
  hline{4}={1-2}{.4pt,solid},
  hline{5}={1-2}{.4pt,solid},
  hline{5}={3}{.4pt,solid},
  vline{3}={4}{.4pt,solid},
  vline{4}={4}{.4pt,solid}}{\toprule
Name&City&Height\\
John&\flmCellsSet{red}Berlin&1.80\\
Anna&\flmCellsSet{green}Zurich&\flmCellsSet{blue}1.65\\
\flmCellsSet{yellow}Two columns by name&&Last\\
&To the end of the table&\\
\bottomrule
}

\begin{flmFloat}{table}{NumCap}
  \flmCellsTable{ccc}{X[-1]X[-1]X[-1]}{hspan=minimal}{\toprule
&Name&Value\\
&Alpha&1\\
&Beta&2\\
\bottomrule
}
  
  
\caption{A simple table.}\flmLabels{\label{ref2}}\end{flmFloat}

See \flmRef{ref2} and \flmRef[S]{ref2}.

\begin{itemize}
\item An item holding a table:
\flmCellsTable{cc}{X[-1]X[-1]}{hspan=minimal,
  cell{1}{1}={}{m, font={\flmCellsHeaderFont}},
  cell{1}{2}={}{m, font={\flmCellsHeaderFont}},
  hline{2}={1}{.4pt,solid},
  hline{2}={2}{.4pt,solid}}{\toprule
Key&Value\\
a&1\\
\bottomrule
}
\end{itemize}

A footnote holding blocks.\flmXNfootnote{Its first paragraph.

\flmCellsTable{cc}{X[-1]X[-1]}{hspan=minimal}{\toprule
c&d\\
\bottomrule
}}

\flmCellsTable{cc}{X[-1]X[-1]}{hspan=minimal,
  cell{1}{1}={}{m, font={\flmCellsHeaderFont}},
  cell{1}{2}={}{m, font={\flmCellsHeaderFont}},
  hline{2}={1}{.4pt,solid},
  hline{2}={2}{.4pt,solid}}{\toprule
Content&Nested\\
A formula \(x^2\), a reference to \flmRef{ref1} and a note\flmXNfootnote{In a cell.}.

  A second paragraph.&\flmCellsTable{cc}{X[-1]X[-1]}{hspan=minimal,
  cell{1}{1}={}{m, font={\flmCellsHeaderFont}},
  cell{1}{2}={}{m, font={\flmCellsHeaderFont}},
  hline{2}={1}{.4pt,solid},
  hline{2}={2}{.4pt,solid}}{\toprule
n&m\\
1&2\\
\bottomrule
}\\
\bottomrule
}

\flmHeading[toc=true]
\section{Mistakes}{\flmLabels{}}

Content outside a cell: 
% flm: dropped content outside \cell
\flmCellsTable{cc}{X[-1]X[-1]}{hspan=minimal}{\toprule
A&B\\
\bottomrule
}

A placement that cannot be read: \flmCellsTable{cc}{X[-1]X[-1]}{hspan=minimal}{\toprule
A&P\\
\bottomrule
}

A command among the styles: \flmCellsTable{c}{X[-1]}{hspan=minimal}{\toprule
S\\
\bottomrule
}

An unknown style: \flmCellsTable{c}{X[-1]}{hspan=minimal,
  cell{1}{1}={}{r}}{\toprule
U\\
\bottomrule
}

A value a style does not take: \flmCellsTable{c}{X[-1]}{hspan=minimal}{\toprule
V\\
\bottomrule
}

Names that cannot be: \flmCellsTable{c}{X[-1]}{hspan=minimal}{\toprule
N\\
\bottomrule
}

Two cells on one position: 
% flm: dropped cell overlapping another
\flmCellsTable{c}{X[-1]}{hspan=minimal}{\toprule
O\\
\bottomrule
}

A table too large: 
% flm: dropped cell making the table too large
\flmCellsTable{c}{X[-1]}{hspan=minimal}{\toprule
T\\
\bottomrule
}

An empty table: 

A blank line inside a row's data: \flmCellsTable{cc}{X[-1]X[-1]}{hspan=minimal}{\toprule
A

B&C\\
\bottomrule
}

Commands outside their table:
\flmCellsTable{cc}{X[-1]X[-1]}{hspan=minimal}{\toprule
x&y 
% flm: dropped & out of place
 z\\
\emph{v 
% flm: dropped \\ out of place
 w}&\\
\bottomrule
}%%
% ===== latex-recomposed:style flmrecomposed.6aba9645.sty =====
\NeedsTeXFormat{LaTeX2e}
\ProvidesPackage{flmrecomposed.6aba9645}
% --- packages ---
% flm.latex.pkg.caption
\RequirePackage{caption}
% flm.latex.pkg.tabularray
\RequirePackage{tabularray}
% flm.latex.pkg.xcolor
\RequirePackage{xcolor}
% flm.latex.pkg.zref-clever
\RequirePackage{zref-clever}
% --- commands and environments the library defines for the document ---
% flm.latex.def.xnotes.footnote
% \flmXNfootnote[<labels>]{<note>}: a note of the kind `footnote' in the FLM document;
% by default it passes the note, its labels first, to \footnote.
% Redefine it to set these notes another way.
\ProvideDocumentCommand\flmXNfootnote{+O{} +m}{\footnote{#1#2}}
% flm.latex.def.floats
% Floats: a float of FLM's figure, table or a configured float type is
% written as \begin{flmFloat}{<type>}{<case>} <content> \caption{<caption>}
% <labels> \end{flmFloat}, <case> being NumCap (numbered, with a caption),
% NumOnly (numbered, no caption: an empty \caption{} shows the number),
% CapOnly (a caption, no number) or Bare (neither).
% \flmFloatSetUseConfig<case>{<type>}{<placement>}: how a float of the case is
%   set. By default a numbered float is set in the environment <type> at
%   <placement>, and an unnumbered one where it is written, in a center
%   environment in which \caption works as in the float <type>; a CapOnly
%   float's caption is set with no number (\caption* of the caption package).
% \flmFloatPlacementArgs<case>: the placement of a float of the case, [tbph]
%   for the numbered cases and [h] for the others.
% The content is centered. Redefine these commands to set floats differently.
\providecommand\flmFloatPlacementArgsNumCap{[tbph]}
\providecommand\flmFloatPlacementArgsNumOnly{[tbph]}
\providecommand\flmFloatPlacementArgsCapOnly{[h]}
\providecommand\flmFloatPlacementArgsBare{[h]}
\def\flm@float@infloat#1#2{\edef\flm@float@begin{\noexpand\begin{#1}#2}\def\flm@float@end{\end{#1}}}
\def\flm@float@incenter#1{\def\flm@float@begin{\begin{center}\def\@captype{#1}}\def\flm@float@end{\end{center}}}
\def\flm@float@starcaption{\let\flm@float@caption\caption\def\caption{\flm@float@caption*}}
\ProvideDocumentCommand\flmFloatSetUseConfigNumCap{m m}{\flm@float@infloat{#1}{#2}}
\ProvideDocumentCommand\flmFloatSetUseConfigNumOnly{m m}{\flm@float@infloat{#1}{#2}}
\ProvideDocumentCommand\flmFloatSetUseConfigCapOnly{m m}{\flm@float@incenter{#1}\flm@float@starcaption}
\ProvideDocumentCommand\flmFloatSetUseConfigBare{m m}{\flm@float@incenter{#1}}
\ProvideDocumentEnvironment{flmFloat}{m m}{%
  \edef\flm@float@placement{\csname flmFloatPlacementArgs#2\endcsname}%
  \csname flmFloatSetUseConfig#2\endcsname{#1}{\flm@float@placement}%
  \flm@float@begin\centering}{\flm@float@end}
% flm.latex.def.cells
% Tables of cells. A table is written
% \flmCellsTable{<natural>}{<fitted>}{<options>}{<rows>}: <natural> and
% <fitted> are two column specifications of the tblr environment of the
% tabularray package (<natural> has c for each column, <fitted> has X[-1]
% for each column),
% <options> are the table's other tblr options (the cells' spans, alignment
% and header font, and the lines), and <rows> are the rows, the cells
% separated by & and each row ended by \\.
% By default the table is set at its natural width with <natural>; where it
% is wider than \flmCellsMaxWidth, it is set again at that width with
% <fitted>. It is set between \flmCellsBeginCenter and \flmCellsEndCenter,
% except where LaTeX sets a single line, as in a cell of another table set
% at its natural width, where no paragraph can begin. A header cell is set
% in \flmCellsHeaderFont. Redefine \flmCellsTable to set tables another way,
% or one of the four other commands to change one part of it.
% A cell's custom style <name> is written \flmCellsSet{<name>} at the start
% of the cell's text: a tabularray table command that applies the tabularray
% cell keys held in \flmCellsKeys<name>. \flmCellsStyleKeys{<name>}{<keys>}
% sets those keys unless keys are already set for <name>: define
% \flmCellsKeys<name> before loading this file, or redefine it after, to
% change how a style looks. The shipped styles green, blue, yellow and red
% set the background colors flmCellsColor<name>, which you can define before
% loading this file.
% tabularray's counter library restores LaTeX's counters after the passes in
% which tabularray sets a cell only to measure it.
% The macros with @ in their names are internals, not a customization point.
\UseTblrLibrary{booktabs}
\UseTblrLibrary{counter}
\providecommand\flmCellsHeaderFont{\bfseries}
\providecommand\flmCellsBeginCenter{\begin{center}}
\providecommand\flmCellsEndCenter{\end{center}}
\providecommand\flmCellsMaxWidth{0.96\linewidth}
\ExplSyntaxOn
\long\def\flm@cells@fit#1#2#3#4{%
  \setbox0=\hbox{\begin{tblr}{colspec={#1},#3}#4\end{tblr}}%
  \dim_compare:nNnTF {\wd0} < {\flmCellsMaxWidth}
    {\leavevmode\box0}
    {\begin{tblr}{width=\flmCellsMaxWidth,colspec={#2},#3}#4\end{tblr}}}
\long\def\flm@cells@iflr#1#2{\mode_if_horizontal:TF {\mode_if_inner:TF {#1} {#2}} {#2}}
\def\flm@cells@set#1{\ifcsname flmCellsKeys#1\endcsname
  \exp_args:Nnv \tblr_set_cell:nn {} {flmCellsKeys#1}\fi}
\ExplSyntaxOff
\ProvideDocumentCommand\flmCellsTable{+m +m +m +m}{%
  \flm@cells@iflr
    {\flm@cells@fit{#1}{#2}{#3}{#4}}%
    {\flmCellsBeginCenter\flm@cells@fit{#1}{#2}{#3}{#4}\flmCellsEndCenter}}
\ProvideDocumentCommand\flmCellsStyleKeys{m m}{%
  \expandafter\providecommand\csname flmCellsKeys#1\endcsname{#2}}
\ifcsname flmCellsSet\endcsname\else
  \NewTableCommand\flmCellsSet[1]{\flm@cells@set{#1}}
  \ProvideDocumentCommand\flmCellsSet{m}{}
\fi
\providecolor{flmCellsColorgreen}{RGB}{200,255,200}
\providecolor{flmCellsColorblue}{RGB}{200,220,255}
\providecolor{flmCellsColoryellow}{RGB}{255,255,200}
\providecolor{flmCellsColorred}{RGB}{255,200,200}
\flmCellsStyleKeys{green}{bg=flmCellsColorgreen}
\flmCellsStyleKeys{blue}{bg=flmCellsColorblue}
\flmCellsStyleKeys{yellow}{bg=flmCellsColoryellow}
\flmCellsStyleKeys{red}{bg=flmCellsColorred}
% flm.latex.def.refs
% \flmRef[<options>]{<label>}: a reference that gives the name of its target
% ("Section 2"), set by \zcref of the zref-clever package.
%
% The options are the ones the document wrote, in FLM's own vocabulary. This
% definition translates each of them by its key -- a value written after the
% key is ignored, as in FLM -- and drops every key it has no translation for:
%
%   noprefix    -> noname                 the number or title alone, with no
%                                         name before it ("2", not
%                                         "Section 2")
%   S           -> S                      the first name capitalized and
%                                         unabbreviated, for a reference that
%                                         opens a sentence
%   capital     -> cap=true,abbrev=false  the names capitalized ("Lemma 1")
%   lowercase   -> cap=false,abbrev=false the names in lowercase ("lemma 1")
%   abbreviated -> abbrev=true            the abbreviated names ("Lem. 1")
%                                         where the reference type declares
%                                         them, the full names otherwise
%   anything else -> dropped              the default wording; a variant of
%                                         your own naming has no counterpart
%                                         in zref-clever, which would stop on
%                                         a key it does not know
%
% Redefine \flmRef to set references differently -- another package, another
% wording, your own variants, no link. Your definition sees the options as the
% document wrote them, in FLM's vocabulary, and maps them as it likes.
%
% The translation uses expl3, LaTeX's programming layer, which is part of the
% format and needs no package; its names read only between \ExplSyntaxOn and
% \ExplSyntaxOff. \keyval_parse:NNn reads the option list as LaTeX reads any
% key-value list: it hands a key written alone to the first macro below and a
% key with its value to the second, which drops the value. \str_case:nnF looks
% the key up in the table and appends what \zcref is to be given, or nothing.
% The macros with @ in their names are internals, not a customization point:
% redefine \flmRef, not its parts.
\ExplSyntaxOn
\def\flm@ref@add#1{\tl_put_right:Nn \flm@ref@options {#1,}}
\def\flm@ref@key#1{\str_case:nnF {#1}
  {{noprefix}{\flm@ref@add{noname}}
   {S}{\flm@ref@add{S}}
   {capital}{\flm@ref@add{cap=true,abbrev=false}}
   {lowercase}{\flm@ref@add{cap=false,abbrev=false}}
   {abbreviated}{\flm@ref@add{abbrev=true}}}
  {}}
\def\flm@ref@pair#1#2{\flm@ref@key{#1}}
\def\flm@ref@zcref#1#2{\zcref[#1]{#2}}
\ProvideDocumentCommand\flmRef{O{} m}{%
  \def\flm@ref@options{}%
  \keyval_parse:NNn \flm@ref@key \flm@ref@pair {#1}%
  \exp_args:No \flm@ref@zcref {\flm@ref@options} {#2}}
\ExplSyntaxOff
% flm.latex.def.label-display-text
% \flmLabelsWithDisplayText{<text>}{<labels>}: makes <labels> refer to this
% place, with <text> as what a reference to them prints. It is the only
% command for that: a heading the class did not number uses it, and so does
% anything else that has a name but no number.
%
% It places no anchor of its own. Under hyperref a sectioning command, starred
% or not, already places one at the title; a \label after it points there.
% \refstepcounter would place a second anchor here, below the title, and a
% link would lead below the heading. So the counter is stepped with
% the plain \stepcounter and the two things this needs are set here:
% \@currentlabel, which is what a \label records, and \@currentcounter,
% which is what zref-clever reads to give the label its reference type. The
% group keeps every change local.
%
% <text> is expanded to make \@currentlabel, and \label does not survive that
% expansion; so a \label inside <text> (for example a footnote's) is dropped
% from the label's text only. At its own place in the document it still makes
% its own label.
%
% A caller with no sectioning command before it -- for example a defined term
% in running text -- writes \phantomsection at its own place, so that a link
% has an anchor of its own to lead to.
%
% Redefine \flmLabelsWithDisplayText to change how such labels are made.
\ifcsname c@flm@internal@refcnt\endcsname\else\newcounter{flm@internal@refcnt}\fi
\ProvideDocumentCommand\flmLabelsWithDisplayText{m m}{%
  \begingroup
    \stepcounter{flm@internal@refcnt}%
    {\let\label\@gobble\protected@xdef\flm@lwdt@text{#1}}%
    \let\@currentlabel\flm@lwdt@text
    \def\@currentcounter{flm@internal@refcnt}%
    #2%
  \endgroup}
% flm.latex.def.zref-clever-type
% The zref-clever reference type of the counter \flmLabelsWithDisplayText
% steps, with every name prefix empty, so that a reference to such a label
% prints its text alone ("Introduction", not "Section Introduction").
% Every other label keeps the name its own type gives it. Redefine these
% settings, or the type of the counter, to put something before the text.
\zcsetup{countertype={flm@internal@refcnt=flm--internal--ref}}
\zcRefTypeSetup{flm--internal--ref}{Name-sg = , name-sg = , Name-pl = , name-pl = ,}
% flm.latex.def.labels
% \flmLabels{<\label commands>}: the labels of one place in the document.
% In text it runs them as written. On a line of a display-math
% environment, where amsmath keeps the last label of a line and loses the
% ones before it, it runs the first and records every further one as an alias:
% a line per alias is written to the .aux file after the last page and read on
% the next run, making the alias refer to what the first label refers to
% (for \ref and for \zcref alike). Redefine \flmLabels to change what a
% group of labels does; \def\flmLabels#1{#1} restores plain LaTeX behavior,
% amsmath's "Multiple \label's" message included.
\ProvideDocumentCommand\flmLabels{m}{%
  \ifmmode\expandafter\flm@labels@math\else\expandafter\@firstofone\fi{#1}}
% Reading the group: the first \label as written, the rest as aliases of it.
\def\flm@labels@math#1{\def\flm@lbl@first{}\flm@lbl@scan#1\flm@lbl@stop}
\def\flm@lbl@scan{\@ifnextchar\label\flm@lbl@take\flm@lbl@rest}
\def\flm@lbl@take\label#1{%
  \ifx\flm@lbl@first\@empty
    \def\flm@lbl@first{#1}%
    \label{#1}%
  \else
    \expandafter\flm@recordalias\expandafter{\flm@lbl@first}{#1}%
  \fi
  \flm@lbl@scan}
\def\flm@lbl@rest#1\flm@lbl@stop{#1}
% Collecting the aliases. A name is recorded once however often the line is
% typeset (amsmath measures a line before setting it).
\def\flm@aliaslist{}
\def\flm@recordalias#1#2{%
  \ifcsname flm@alias@#2\endcsname\else
    \expandafter\gdef\csname flm@alias@#2\endcsname{}%
    \g@addto@macro\flm@aliaslist{\flm@alias@entry{#2}{#1}}%
  \fi}
% Writing them, from the hook that runs after the last page has been shipped
% out: every \newlabel of that page has reached the .aux file by then, so the
% alias lines come after them. \AtEndDocument would run too early.
\def\flm@alias@write#1#2{\immediate\write\@auxout{\string\flm@aliaslabel{#1}{#2}}}
\AddToHook{enddocument/afterlastpage}{\let\flm@alias@entry\flm@alias@write\flm@aliaslist}
% Reading them back: the alias refers to what the original refers to. The
% assignments are global because the .aux file is read inside a group.
\def\flm@aliaslabel#1#2{%
  \ifcsname r@#2\endcsname
    \global\expandafter\let\csname r@#1\expandafter\endcsname\csname r@#2\endcsname
  \fi
  \ifcsname Z@R@#2\endcsname
    \global\expandafter\let\csname Z@R@#1\expandafter\endcsname\csname Z@R@#2\endcsname
  \fi}
% flm.latex.def.heading
% \flmHeading[<keys>]<sectioning command><star>{<title>}{<labels>}: a heading
% and the labels that refer to it, written on two lines for readability, as in
%   \flmHeading[toc=true]
%   \section*{Title}{\flmLabels{\label{ref1}}}
% It runs the sectioning command as written, then makes the labels refer to
% the heading: by its number when the class numbered it, by its title when it
% did not (a starred heading, or a level below secnumdepth). It finds out
% which without patching anything: a class numbers a heading through
% \refstepcounter, which records the counter's name in \@currentcounter.
%
% The optional argument is a key-value list. It is written here, on the hook,
% and never on the sectioning command: a starred sectioning command takes no
% optional argument, and giving one to \section* is not an error, it is
% silently wrong -- \section*[Short]{Long} takes one token as its heading,
% so the opening bracket alone becomes the heading and Short, the closing
% bracket and Long all run together into the following paragraph. Two keys
% are defined:
%
%   toc=true|false     whether the heading is listed in the table of contents.
%                      Absent, it is listed as LaTeX lists it. After a starred
%                      sectioning command with toc=true, which LaTeX never
%                      lists, \addcontentsline{toc}{<level>}{<entry>} adds it:
%                      <level> is the command's own name (section for
%                      \section), <entry> the alternative title if alt was
%                      written, else the title; under hyperref the starred
%                      command has already placed the anchor the line links
%                      to. On an unstarred heading toc=false is not honored:
%                      the heading stays listed and a warning says so.
%
%   alt={Short title}  an alternative title: what goes into the table of
%                      contents, the running head and the PDF bookmark, while
%                      the page keeps the full title. It is handed to the
%                      sectioning command as its own optional argument, so it
%                      does what \section[Short title]{Full title} does, and
%                      nothing else: a reference to an unnumbered heading still
%                      prints the full title. A starred heading does not hand
%                      it to the sectioning command; where it is written
%                      toc=true, alt is its line in the table of contents.
%
% What counts is whether the key was written, not what it holds. alt={} is an
% empty alternative title and is handed over as \section[]{Full title}, which
% is how LaTeX asks for an empty table-of-contents entry; no optional argument
% at all, and an empty one ([]), hand nothing over and leave the sectioning
% command exactly as it was. Writing alt with no value is LaTeX's own error,
% as is a key that is not defined here, and a value of toc other than true or
% false. The values and the record that alt was written are cleared at every
% call, so one heading's keys never reach the next. The recomposed LaTeX writes
% toc on every heading and alt where the heading was given an alternative
% title; both keys are also there for a document written by hand and for a
% user who redefines the hook.
%
% Redefine \flmHeading to change what a heading does. The block below uses
% expl3, LaTeX's programming layer, which is part of the format and needs no
% package; its names read only between \ExplSyntaxOn and \ExplSyntaxOff.
% \keys_define:nn declares the keys and \keys_set:nn reads a list of them;
% \bool_if:NTF asks whether the key was written. The three declarations are
% guarded on their own names, as every declaration in this file is, so that two
% of these style files can be loaded in one LaTeX run. The macros with @ in
% their names are internals, not a customization point: redefine \flmHeading,
% not its parts.
\ExplSyntaxOn
\ifcsname flm@heading@alt\endcsname\else\tl_new:N \flm@heading@alt\fi
\ifcsname flm@heading@given\endcsname\else\bool_new:N \flm@heading@given\fi
\ifcsname flm@heading@toc\endcsname\else\tl_new:N \flm@heading@toc\fi
\keys_define:nn {flm/heading}
  {alt .code:n = {\tl_set:Nn \flm@heading@alt {#1}
                  \bool_set_true:N \flm@heading@given} ,
   alt .value_required:n = true ,
   toc .choices:nn = {true, false} {\tl_set_eq:NN \flm@heading@toc \l_keys_choice_tl} ,
   toc .value_required:n = true}
\def\flm@heading@withalt#1#2#3{#2[#1]{#3}}
% The contents line of a starred heading: <level> then <entry>.
\def\flm@heading@addtoc#1#2{\addcontentsline{toc}{#1}{#2}}
% toc=false on an unstarred heading: the heading stays listed.
\def\flm@heading@notoc#1{\PackageWarning{flm}{toc=false~on~the~unstarred~
  \string#1~is~not~honored~by~the~default~\string\flmHeading;~the~heading~
  stays~in~the~table~of~contents}}
\ProvideDocumentCommand\flmHeading{O{} m s m m}{%
  \tl_clear:N \flm@heading@alt
  \bool_set_false:N \flm@heading@given
  \tl_clear:N \flm@heading@toc
  \keys_set:nn {flm/heading} {#1}%
  \def\@currentcounter{}%
  \IfBooleanTF{#3}
    {#2*{#4}%
     \str_if_eq:VnT \flm@heading@toc {true}
       {\edef\flm@heading@level{\cs_to_str:N #2}%
        \bool_if:NTF \flm@heading@given
          {\exp_args:NVV \flm@heading@addtoc \flm@heading@level \flm@heading@alt}
          {\exp_args:NV \flm@heading@addtoc \flm@heading@level {#4}}}}
    {\str_if_eq:VnT \flm@heading@toc {false} {\flm@heading@notoc{#2}}%
     \bool_if:NTF \flm@heading@given
       {\exp_args:NV \flm@heading@withalt \flm@heading@alt #2{#4}}
       {#2{#4}}}%
  \ifx\@currentcounter\@empty
    \flmLabelsWithDisplayText{#4}{#5}%
  \else
    #5%
  \fi}
\ExplSyntaxOff
% --- definitions written from the configuration of FLM features ---
% flm@float@reftype@table (the floats feature's, written from its configuration): declared unless \flm@float@reftype@table is already defined, so a definition made earlier under that name is kept
\ifcsname flm@float@reftype@table\endcsname\else\zcRefTypeSetup{table}{Name-sg={Table},name-sg={Tab.},Name-pl={Tables},name-pl={Tabs.}}\def\flm@float@reftype@table{}\fi
\endinput
